\chapter*{第一版后记}
\addcontentsline{toc}{chapter}{第一版后记}
\markboth{第一版后记}{第一版后记}

本书到第十二章结束。

从一份二十页的社团讲义到一部十二章的教材，中间隔着的不是篇幅，
而是三件事：把“章节”变成“章—节—目”，把“说法”从“事实”里
分出来，把“感想”变成可检验的假设。这三件事里没有一件
是奶蛙本身的“故事”，它们都是\keyw{关于如何研究奶蛙}的
元工作。如果本书对读者最有用的部分是这些，那它就没白写。

有几句话想在最后再说一遍。

第一，\textbf{本书没有任何一项来自作者亲自采集的经验数据}。
第十一章的问卷、示例数据表，都是供修订者与后续研究者使用的
\keyw{设计}与\keyw{教学样例}，不是本书的测量结果。
谁若把本书当作“已有结论”来引用，请回到第七章的假设
与第十二章的可复现性清单，先看它的证伪条件。

第二，\textbf{本书对具体个人保持沉默}。
奶蛙起源考据指向的创作者，以及一切出现在社群材料里的
真实个体，本书只作指称，不作推断——尤其不转述任何关于
其健康、生死、私德的内容。这是证据分级制度里
最难执行、也最不能妥协的一条。

第三，\textbf{本书把“乐子”当作研究对象，但不消费它}。
我们承认奶蛙二创中存在未经授权的形象再生产，
这在版权法上有争议；本书讨论它，是为了理解
“正典—变种”结构如何运作，而不是为了鼓励或背书
任何一种具体的再生产行为。

欢迎各高校同好会以本校的奶蛙活动为基础，
对本书进行修订与续编——更新事实手册中的条目与访问日期，
补充本校的案例与数据，重做第十一章的问卷。
修订时请遵守 CC BY-NC-SA 4.0 协议的三项义务，
并在新的“事实手册”里保留本书的证据分级制度。

如果读完这本书，你能同时做到两件事——
\keyw{看懂一只青蛙为什么让人捧腹大笑}，
并且\keyw{说不清哪一句是事实、哪一句是说法}——
那么它可能没有完全失败。

\vspace{2em}
\begin{flushright}
\BookOrg\\
\BookDate
\end{flushright}
